\documentclass[11pt,a4paper]{article}
% Common preamble for the hint sheets (\input from each file)
\usepackage[utf8]{inputenc}
\usepackage[T1]{fontenc}
\usepackage{lmodern}
\usepackage{amsmath,amssymb,amsthm,mathtools}
\usepackage[a4paper,margin=2.3cm]{geometry}
\usepackage{enumitem}
\usepackage{xcolor}
\usepackage[hidelinks]{hyperref}
\renewcommand{\leq}{\leqslant}
\renewcommand{\le}{\leqslant}
\renewcommand{\geq}{\geqslant}
\renewcommand{\ge}{\geqslant}
\setlength{\parindent}{0pt}
\setlength{\parskip}{0.35em}
% \hintstitle{Sheet title}{Class work / Homework}
\newcommand{\hintstitle}[2]{\begin{center}{\LARGE\bfseries Hints: #1\par}\vspace{0.5em}{\large #2\par}\end{center}
\noindent\textit{The hints to each problem are cumulative. Read one, try again, and open the next only if you are still stuck.}\par\medskip\hrule\medskip}
\newcommand{\problem}[1]{\par\bigskip{\large\bfseries Problem #1}\par\nopagebreak\smallskip}
\newcommand{\subpart}[1]{\par\smallskip\textbf{(#1)}\par\nopagebreak}
\newcommand{\hint}[1]{\par\smallskip\noindent\textcolor{black!70}{\textbf{Hint #1.}}\enspace}
\newcommand{\note}[1]{\par\smallskip\noindent\textit{Note.} #1\par}
\providecommand{\lt}{<}
\providecommand{\gt}{>}
\providecommand{\tightlist}{\setlength{\itemsep}{0pt}\setlength{\parskip}{0pt}}
\emergencystretch=3em

\begin{document}
\hintstitle{Linear algebra -- 6. Linear algebra in combinatorics}{Class work}

\problem{1}
\hint{1}
Oddtown: over \(\mathbb F_2\) the Gram matrix \((\langle 1_{A_i},1_{A_j}\rangle)\) is the identity. So the vectors \(1_{A_i}\in\mathbb F_2^n\) are independent. The singletons show sharpness.
\hint{2}
Eventown: now every inner product, including \(\langle v,v\rangle\), is \(0\). So \(W=\operatorname{span}\{1_{A_i}\}\) satisfies \(W\subseteq W^\perp\), and \(\dim W+\dim W^\perp=n\) gives \(\dim W\le\lfloor n/2\rfloor\). The family lies inside \(W\), so \(m\le|W|\). Sharp: all unions of the pairs \(\{1,2\},\{3,4\},\dots\).
\hint{3}
The difference: in Oddtown the vectors are independent, so you count dimensions. Eventown families can be closed under symmetric difference, i.e.\ be a whole subspace, so you count the elements of a subspace.

\problem{2}
\hint{1}
The answer is \(2n\) (take \(\pm e_i\)). For the bound, put \(M=\sum_i v_iv_i^T\) and compute \(\operatorname{tr}M\), \(\operatorname{tr}M^2\), \(\operatorname{tr}M^3\) in terms of the inner products.
\hint{2}
Let \(S=\sum_{i\ne j}\langle v_i,v_j\rangle^2\). Then \(\operatorname{tr}M=m\) and \(\operatorname{tr}M^2=m+S\). In \(\operatorname{tr}M^3=\sum\langle v_i,v_j\rangle\langle v_j,v_k\rangle\langle v_k,v_i\rangle\), terms with \(i,j,k\) distinct vanish by hypothesis, so \(\operatorname{tr}M^3=m+3S\).
\hint{3}
Apply Cauchy--Schwarz to the eigenvalues \(\mu_l\ge0\) of \(M\). First \((\sum\mu^2)^2\le\sum\mu\cdot\sum\mu^3\) gives \(S\le m\). Then \((\sum\mu)^2\le n\sum\mu^2\) gives \(m^2\le n(m+S)\le2nm\).

\problem{3}
\hint{1}
Small cases: the minimum is \(0\) for \(n=1\) and \(2\) for \(n=2\). For \(n\ge3\), try to write \(a_{ij}\) as a sum of few products \(f(i)g(j)\) vanishing on the diagonal: \((i-j)^2=i^2-2ij+j^2\) has rank \(\le3\).
\hint{2}
Rank \(2\) is impossible for \(n\ge3\). The first two rows are independent, so row 3 \(=x\cdot\)row 1\(+y\cdot\)row 2. Columns 1 and 2 force \(y>0\) and \(x>0\), and then column 3 gives \(0=xa_{13}+ya_{23}>0\). Answer: \(3\) for \(n\ge3\).

\problem{4}
\hint{1}
Independence of the ordering: reordering conjugates \(M\) by a permutation matrix. Then set up an induction on \(n\). Order the subsets as: non-empty subsets of \([n-1]\), then \(\{n\}\), then \(S\cup\{n\}\) for non-empty \(S\subseteq[n-1]\).
\hint{2}
In this order \(M_n=\begin{pmatrix}M_{n-1}&0&M_{n-1}\\0&1&\mathbf 1^T\\M_{n-1}&\mathbf 1&J\end{pmatrix}\). Subtract the first block row from the third, the first block column from the third, and then the middle row from the last block.
\hint{3}
You get \(\det M_n=\det M_{n-1}\cdot\det(-M_{n-1})=-(\det M_{n-1})^2\), since \(M_{n-1}\) has odd size. With \(\det M_1=1\): \(\det M_n=-1\) for \(n\ge2\).

\problem{5}
\hint{1}
The affine dimension of a point set equals the rank of its centred Gram matrix \(-\frac12HDH\), where \(D\) is the matrix of squared distances and \(H=E-\frac1{2n}J\). Here \(D=b^2(J-E)+(a^2-b^2)P\), with \(P\) the involution swapping \(x_i\leftrightarrow y_i\).
\hint{2}
\(P\) commutes with \(J\) and \(H\). Split \(\mathbb R^{2n}\) into vectors symmetric under \(P\) (dimension \(n\), containing \(\mathbf 1\)) and antisymmetric ones (dimension \(n\)). On \(\mathbf 1^\perp\) the Gram matrix acts as \(b^2-\frac{a^2}{2}\) on the symmetric part (dimension \(n-1\)) and as \(\frac{a^2}{2}\) on the antisymmetric part.
\hint{3}
So the rank is \(n\) if \(b^2=a^2/2\) and \(2n-1\) otherwise. But the points span at most \(d\le2n-2\) dimensions. Hence \(b^2=a^2/2\) and the points lie in an \(n\)-dimensional affine subspace.
\end{document}
